\section{Marco teórico}
\label{sec:marco}

El presente marco conceptual establece el piso teórico interdisciplinario entre
la inteligencia artificial aplicada, la teoría clásica de control estocástico y
la informática industrial moderna. Se organiza en tres ejes conceptuales, y cada
uno se vincula explícitamente con el objetivo específico que sustenta.

\subsection{Eje 1: Gemelos digitales, sombras digitales y la exigencia de
consistencia}
\label{sec:marco-eje1}

\emph{Este eje sustenta la formulación del problema y, de manera directa, la
delimitación del alcance declarada en OE3: el módulo se concibe como componente
arquitectónico y no como herramienta de control.}

El paradigma de los gemelos digitales ha transitado desde la conceptualización
aeroespacial hacia arquitecturas troncales de la Industria 4.0.
\textcite{fuller2020} postulan que un gemelo digital trasciende la simulación
tridimensional estática al requerir una sincronización continua con el activo
físico. En la industria pesada, \textcite{tao2019} y \textcite{liu2021}
argumentan que la principal promesa de esta tecnología es habilitar el
mantenimiento predictivo y la toma de decisiones autónoma mediante inteligencia
artificial, articulando la visión de sistemas ciberfísicos originalmente
planteada por \textcite{lee2008}.

Sin embargo, al analizar críticamente estos enfoques sistémicos, se evidencia
que asumen tácitamente la existencia de una red de telemetría libre de ruido.
\textcite{sepasgozar2021} aporta la distinción conceptual decisiva: una réplica
que se limita a recibir y espejar datos crudos constituye una \emph{sombra
digital}, mientras que un verdadero \emph{gemelo digital} procesa y depura la
incertidumbre subyacente. Esta distinción no es semántica sino arquitectónica, y
determina dónde debe ubicarse la responsabilidad de la calidad del dato dentro
del sistema.

De esta discusión se deriva la posición teórica del proyecto: si la depuración
de la incertidumbre es condición constitutiva del gemelo digital, entonces el
estimador estocástico no es un accesorio de preprocesamiento sino un componente
de primer orden de la arquitectura, y debe diseñarse, documentarse y validarse
como tal.

\subsection{Eje 2: Observabilidad parcial y degradación de la telemetría
industrial}
\label{sec:marco-eje2}

\emph{Este eje fundamenta OE1, la caracterización estadística de las
perturbaciones, y OE2, la formulación del modelo dinámico con estados no
sensados.}

La aplicación de los conceptos anteriores al chancado primario choca con un
desafío físico inherente: la severidad del entorno. La condición operativa de
una estación de chancado —vibración mecánica continua de gran magnitud y alta
concentración de polvo abrasivo en suspensión— degrada de forma sistemática las
señales de los transductores acústicos de nivel.

El resultado es una situación de \textbf{observabilidad parcial}, en la que el
estado real del sistema difiere sustancialmente del dato reportado en consola.
Esta condición tiene dos manifestaciones distintas que conviene separar. La
primera es la \emph{corrupción por ruido} de las variables que sí se miden: el
volumen es observado, pero con una varianza que compromete su utilidad para la
automatización. La segunda, más severa, es la \emph{ausencia total de medición}
de variables físicamente relevantes: la tasa instantánea de extracción del
chancador no cuenta con instrumentación directa y solo puede inferirse.

Las revisiones de \textcite{qu2023} y \textcite{ebad2026} confirman que la
integración y la calidad de datos permanecen como desafíos abiertos en la
adopción de gemelos digitales mineros, y que el esfuerzo investigativo se ha
concentrado en las capas superiores de predicción antes que en esta capa
intermedia de consistencia.

Frente a este diagnóstico, el proyecto adopta una posición metodológica
específica: la solución no radica en la infraestructura de hardware —camino que
implicaría reemplazo de instrumentación o adición de redundancia sensorial—
sino en filtros algorítmicos operando sobre la señal existente. Esta elección es
la que hace del problema un problema de ingeniería informática.

\subsection{Eje 3: Estimación óptima de estados como capa de consistencia}
\label{sec:marco-eje3}

\emph{Este eje sustenta OE2, OE3 y OE4: aporta las ecuaciones del modelo, el
algoritmo a implementar y el criterio estadístico bajo el cual se valida.}

Para resolver la observabilidad parcial y reconstruir matemáticamente el volumen
de mineral, el proyecto se fundamenta en las ecuaciones estocásticas lineales de
espacio de estados en tiempo discreto:

\begin{align}
\mathbf{x}_{k+1} &= \mathbf{A}\mathbf{x}_{k} + \mathbf{B}\mathbf{u}_{k}
                    + \mathbf{w}_{k}, & \mathbf{w}_{k} &\sim
                    \mathcal{N}(\mathbf{0}, \mathbf{Q}) \label{eq:estado}\\
\mathbf{y}_{k}   &= \mathbf{H}\mathbf{x}_{k} + \mathbf{v}_{k}, &
                    \mathbf{v}_{k} &\sim \mathcal{N}(\mathbf{0}, \mathbf{R})
                    \label{eq:medicion}
\end{align}

donde $\mathbf{x}_{k} \in \mathbb{R}^{n}$ representa el vector de estados
latentes reales del activo, $\mathbf{u}_{k} \in \mathbb{R}^{m}$ las entradas
conocidas registradas por el SCADA, y $\mathbf{y}_{k} \in \mathbb{R}^{p}$ el
vector de observaciones afectadas por las condiciones ambientales. Las
perturbaciones internas del proceso y el ruido instrumental se modelan mediante
vectores gaussianos independientes de media cero, cuyas matrices de covarianza
$\mathbf{Q}$ y $\mathbf{R}$ delimitan el comportamiento estadístico del entorno.

La teoría de estimación óptima establece que el Filtro de Kalman es el estimador
recursivo de varianza mínima para sistemas lineales inmersos en ruido gaussiano
blanco \parencite{kalman1960,simon2006}. Su operación se organiza en un ciclo
predictor-corrector: la etapa de predicción propaga el estado y su incertidumbre
mediante el modelo físico, y la etapa de corrección fusiona esa predicción con
la medición ruidosa ponderando ambas fuentes según su confiabilidad relativa.

\begin{align}
\hat{\mathbf{x}}_{k|k-1} &= \mathbf{A}\hat{\mathbf{x}}_{k-1|k-1}
                            + \mathbf{B}\mathbf{u}_{k} \label{eq:pred-x}\\
\mathbf{P}_{k|k-1} &= \mathbf{A}\mathbf{P}_{k-1|k-1}\mathbf{A}^{\mathsf{T}}
                      + \mathbf{Q} \label{eq:pred-P}\\
\mathbf{K}_{k} &= \mathbf{P}_{k|k-1}\mathbf{H}^{\mathsf{T}}
                  \left(\mathbf{H}\mathbf{P}_{k|k-1}\mathbf{H}^{\mathsf{T}}
                  + \mathbf{R}\right)^{-1} \label{eq:ganancia}\\
\hat{\mathbf{x}}_{k|k} &= \hat{\mathbf{x}}_{k|k-1}
                          + \mathbf{K}_{k}\left(\mathbf{y}_{k}
                          - \mathbf{H}\hat{\mathbf{x}}_{k|k-1}\right)
                          \label{eq:corr-x}
\end{align}

La ganancia $\mathbf{K}_{k}$ de la ecuación~\eqref{eq:ganancia} concentra la
lógica del estimador: cuando la incertidumbre del modelo $\mathbf{P}_{k|k-1}$ es
grande respecto de la del sensor $\mathbf{R}$, la ganancia crece y el filtro
confía en la medición; en el caso inverso, la ganancia decrece y el filtro
confía en su propia predicción física. Esta ponderación es adaptativa y se
recalcula en cada ciclo.

La pertinencia de este algoritmo para el proyecto radica en tres propiedades. La
primera es su \textbf{eficiencia computacional nativa}: el ciclo se reduce a
operaciones de álgebra lineal sobre matrices de baja dimensión, lo que permite
su ejecución con costo despreciable, requisito fundamental para un gemelo digital
en tiempo real. La segunda es su \textbf{fundamento físico}: a diferencia de un
suavizador estadístico genérico, el Filtro de Kalman incorpora las leyes de
conservación del proceso en la matriz $\mathbf{A}$, de modo que la estimación
respeta la física del activo. La tercera es su \textbf{capacidad de reconstruir
estados no medidos}, siempre que se satisfaga la condición de observabilidad del
par $(\mathbf{A}, \mathbf{H})$, propiedad que se verifica explícitamente en la
sección~\ref{sec:unidad}.

La viabilidad de esta aproximación en el dominio minero cuenta con respaldo
empírico. \textcite{apelt2002} demostraron la medición inferencial de parámetros
internos de molinos semiautógenos mediante estimación de estados, y
\textcite{barzaga2025} aplicaron filtros de Kalman extendido y \emph{unscented}
a la estimación de variables no observables en lixiviación en pilas de cobre.
Ambos trabajos acreditan que la reconstrucción de estados latentes es efectiva
en procesos de tratamiento de minerales; la contribución del presente proyecto
consiste en formular esa capacidad como componente arquitectónico de un gemelo
digital, con alcance delimitado y validación estadística formal.
